\chapter{Calculating by hand}\label{ch:byhand}

Each test in this guide, done with a pencil on a handful of numbers. Every section follows the same order of thinking: what we want to measure, what kind of data it is, how such data behave, which test fits and why, and then the arithmetic, one step under another, with nothing skipped. The typed-out working below each section is produced by the notebook \texttt{by\_hand.ipynb}, which computes every line and checks the final answer against \texttt{scipy.stats}; change its numbers and every line changes with them. Work each example on paper first, then compare.

\section{How to use this chapter}
\input{generated/nb_by_hand_1}

\section{A coin: the binomial test}
\input{generated/nb_by_hand_2}
\subsection*{The working, typeset}
\begin{align*}
P(X = k) &= \binom{n}{k}\,\frac{1}{2^{n}} \qquad \text{\small (chance of exactly k heads)} \\
2^{10} &= 1024 \qquad \text{\small (number of possible sequences)} \\
\binom{10}{8} &= \frac{10!}{8!\,2!} = 45 \\
\binom{10}{9} &= \frac{10!}{9!\,1!} = 10 \\
\binom{10}{10} &= \frac{10!}{10!\,0!} = 1 \\
P(X \ge 8) &= \frac{45 + 10 + 1}{1024} \\
 &= \frac{56}{1024} \\
 &= 0.0547 \qquad \text{\small (one tail)} \\
p &= 2 \times 0.0547 \qquad \text{\small (the other tail, 0 to 2 heads, is the same size)} \\
 &= 0.1094
\end{align*}


\section{The z-test for a mean}
\input{generated/nb_by_hand_3}
\subsection*{The working, typeset}
\begin{align*}
\bar{x} &= \frac{20.3 + 20.1 + 20.4 + 20.2}{4} \\
 &= \frac{81}{4} \\
 &= 20.25 \\
\mathrm{SE} &= \frac{\sigma}{\sqrt{n}} \\
 &= \frac{0.2}{\sqrt{4}} \\
 &= \frac{0.2}{2} \\
 &= 0.1 \\
z &= \frac{\bar{x} - \mu_0}{\mathrm{SE}} \\
 &= \frac{20.25 - 20}{0.1} \\
 &= \frac{0.25}{0.1} \\
 &= 2.5 \\
p &= 2 \times P(Z > 2.5) \qquad \text{\small (two-sided: too high or too low)} \\
 &= 2 \times 0.0062 \\
 &= 0.0124
\end{align*}


\section{The t-test for a trend}
\input{generated/nb_by_hand_4}
\subsection*{The working, typeset}
\begin{align*}
\bar{x} &= \frac{1 + 2 + 3 + 4 + 5}{5} \\
 &= 3 \\
\bar{y} &= \frac{0.1 + 0.25 + 0.2 + 0.4 + 0.45}{5} \\
 &= \frac{1.4}{5} \\
 &= 0.28 \\
S_{xx} &= \sum (x_i - \bar{x})^2 \\
 &= (-2)^2 + (-1)^2 + (0)^2 + (1)^2 + (2)^2 \\
 &= 4 + 1 + 0 + 1 + 4 \\
 &= 10 \\
S_{xy} &= \sum (x_i - \bar{x})(y_i - \bar{y}) \\
 &= (-2)(-0.18) + (-1)(-0.03) + (0)(-0.08) + (1)(0.12) + (2)(0.17) \\
 &= (0.36) + (0.03) + (0) + (0.12) + (0.34) \\
 &= 0.85 \\
b &= \frac{S_{xy}}{S_{xx}} = \frac{0.85}{10} \\
 &= 0.085 \qquad \text{\small (degrees C per year)} \\
a &= \bar{y} - b\,\bar{x} = 0.28 - 0.085 \times 3 \\
 &= 0.025 \\
\hat{y}_i &= a + b\,x_i = 0.11,\ 0.195,\ 0.28,\ 0.365,\ 0.45 \\
e_i &= y_i - \hat{y}_i = -0.01,\ 0.055,\ -0.08,\ 0.035,\ 0 \\
\mathrm{SSE} &= \sum e_i^2 = (-0.01)^2 + (0.055)^2 + (-0.08)^2 + (0.035)^2 + (0)^2 \\
 &= 0.01075 \\
s^2 &= \frac{\mathrm{SSE}}{n - 2} = \frac{0.01075}{3} \\
 &= 0.003583 \\
\mathrm{SE}(b) &= \sqrt{\frac{s^2}{S_{xx}}} = \sqrt{\frac{0.003583}{10}} \\
 &= \sqrt{0.0003583} \\
 &= 0.01893 \\
t &= \frac{b}{\mathrm{SE}(b)} = \frac{0.085}{0.01893} \\
 &= 4.49 \\
\mathrm{df} &= n - 2 = 3 \\
p &= 2 \times P(T_3 > 4.49) = 0.0206 \qquad \text{\small (from a t table or software)}
\end{align*}


\section{Welch's t-test: two groups}
\input{generated/nb_by_hand_5}
\subsection*{The working, typeset}
\begin{align*}
\bar{x}_1 &= \frac{82 + 90 + 78 + 94}{4} = \frac{344}{4} \\
 &= 86 \\
s_1^2 &= \frac{\sum (x - \bar{x}_1)^2}{n_1 - 1} \\
 &= \frac{(-4)^2 + (4)^2 + (-8)^2 + (8)^2}{3} \\
 &= \frac{16 + 16 + 64 + 64}{3} \\
 &= \frac{160}{3} = 53.3333 \\
\bar{x}_2 &= \frac{105 + 112 + 98 + 120 + 110}{5} = \frac{545}{5} \\
 &= 109 \\
s_2^2 &= \frac{\sum (x - \bar{x}_2)^2}{n_2 - 1} \\
 &= \frac{(-4)^2 + (3)^2 + (-11)^2 + (11)^2 + (1)^2}{4} \\
 &= \frac{16 + 9 + 121 + 121 + 1}{4} \\
 &= \frac{268}{4} = 67 \\
\mathrm{SE} &= \sqrt{\frac{s_1^2}{n_1} + \frac{s_2^2}{n_2}} \\
 &= \sqrt{\frac{53.3333}{4} + \frac{67}{5}} \\
 &= \sqrt{13.3333 + 13.4} \\
 &= \sqrt{26.7333} = 5.1704 \\
t &= \frac{\bar{x}_1 - \bar{x}_2}{\mathrm{SE}} = \frac{86 - 109}{5.1704} \\
 &= \frac{-23}{5.1704} = -4.448 \\
\mathrm{df} &= \frac{\left(\frac{s_1^2}{n_1} + \frac{s_2^2}{n_2}\right)^2}{\frac{(s_1^2/n_1)^2}{n_1 - 1} + \frac{(s_2^2/n_2)^2}{n_2 - 1}} \\
 &= \frac{(26.7333)^2}{\frac{13.3333^2}{3} + \frac{13.4^2}{4}} \\
 &= \frac{714.6711}{59.2593 + 44.89} \\
 &= 6.86 \\
p &= 2 \times P(T_{6.86} > 4.448) = 0.00313
\end{align*}


\section{Mann--Whitney U: two groups by ranks}
\input{generated/nb_by_hand_6}
\subsection*{The working, typeset}
\begin{align*}
\text{sorted} &= 78_{A},\ 82_{A},\ 90_{A},\ 94_{A},\ 98_{B},\ 105_{B},\ 110_{B},\ 112_{B},\ 120_{B} \\
\text{ranks of A} &= 2,\ 3,\ 1,\ 4 \\
R_A &= 2 + 3 + 1 + 4 = 10 \\
U_A &= R_A - \frac{n_A(n_A + 1)}{2} \\
 &= 10 - \frac{4 \times 5}{2} \\
 &= 10 - 10 = 0 \\
U_B &= n_A n_B - U_A = 4 \times 5 - 0 = 20 \\
U &= \min(U_A, U_B) = 0 \\
\text{ways to choose A's ranks} &= \binom{9}{4} = 126 \\
P(U \le 0) &= \frac{1}{126} \qquad \text{\small (only one arrangement puts all of A at the bottom)} \\
p &= 2 \times \frac{1}{126} = 0.0159 \qquad \text{\small (two-sided)}
\end{align*}


\section{Correlation: Pearson and Spearman}
\input{generated/nb_by_hand_7}
\subsection*{The working, typeset}
\begin{align*}
\bar{x} &= 0.16, \quad \bar{y} = 98.8 \\
x_i - \bar{x} &= -1.36,\ -0.56,\ -0.06,\ 0.64,\ 1.34 \\
y_i - \bar{y} &= 13.2,\ 5.2,\ 2.2,\ -6.8,\ -13.8 \\
S_{xx} &= 1.8496 + 0.3136 + 0.0036 + 0.4096 + 1.7956 = 4.372 \\
S_{yy} &= 174.24 + 27.04 + 4.84 + 46.24 + 190.44 = 442.8 \\
S_{xy} &= (-17.952) + (-2.912) + (-0.132) + (-4.352) + (-18.492) = -43.84 \\
r &= \frac{S_{xy}}{\sqrt{S_{xx}\,S_{yy}}} = \frac{-43.84}{\sqrt{4.372 \times 442.8}} \\
 &= \frac{-43.84}{\sqrt{1935.9216}} = \frac{-43.84}{43.9991} \\
 &= -0.9964 \\
t &= \frac{r\sqrt{n - 2}}{\sqrt{1 - r^2}} = \frac{-0.9964 \times \sqrt{3}}{\sqrt{1 - 0.9928}} \\
 &= \frac{-1.7258}{0.085} = -20.31 \\
p &= 2 \times P(T_3 > 20.31) = 0.00026 \\
\text{ranks of } x &= 1,\ 2,\ 3,\ 4,\ 5 \\
\text{ranks of } y &= 5,\ 4,\ 3,\ 2,\ 1 \\
d_i &= -4,\ -2,\ 0,\ 2,\ 4 \\
\rho &= 1 - \frac{6\sum d_i^2}{n(n^2 - 1)} = 1 - \frac{6 \times 40}{5 \times 24} \\
 &= -1
\end{align*}


\section{Mann--Kendall: a trend by counting}
\input{generated/nb_by_hand_8}
\subsection*{The working, typeset}
\begin{align*}
\text{pairs} &= \frac{n(n-1)}{2} = \frac{6 \times 5}{2} = 15 \\
\text{year 1 vs later} &= -1,\ -1,\ -1,\ -1,\ -1 \quad (sum -5) \\
\text{year 2 vs later} &= +1,\ -1,\ -1,\ -1 \quad (sum -2) \\
\text{year 3 vs later} &= -1,\ -1,\ -1 \quad (sum -3) \\
\text{year 4 vs later} &= -1,\ -1 \quad (sum -2) \\
\text{year 5 vs later} &= +1 \quad (sum 1) \\
S &= -11 \qquad \text{\small (more falls than rises)} \\
\tau &= \frac{S}{n(n-1)/2} = \frac{-11}{15} = -0.7333 \\
\mathrm{Var}(S) &= \frac{n(n-1)(2n+5)}{18} = \frac{6 \times 5 \times 17}{18} \\
 &= 28.3333 \\
Z &= \frac{S + 1}{\sqrt{\mathrm{Var}(S)}} = \frac{-10}{\sqrt{28.3333}} = \frac{-10}{5.3229} \qquad \text{\small (+1: continuity correction)} \\
 &= -1.879 \\
p &= 2 \times P(Z > 1.879) = 0.0603
\end{align*}


\section{Chi-square: counts in a table}
\input{generated/nb_by_hand_9}
\subsection*{The working, typeset}
\begin{align*}
\text{row totals} &= 60,\ 40; \quad \text{column totals}\ 52,\ 48; \quad N = 100 \\
E_{11} &= \frac{60 \times 52}{100} = 31.2 \qquad \text{\small (before no, dark no)} \\
E_{12} &= \frac{60 \times 48}{100} = 28.8 \qquad \text{\small (before no, dark yes)} \\
E_{21} &= \frac{40 \times 52}{100} = 20.8 \qquad \text{\small (before yes, dark no)} \\
E_{22} &= \frac{40 \times 48}{100} = 19.2 \qquad \text{\small (before yes, dark yes)} \\
\frac{(O - E)^2}{E}_{11} &= \frac{(42 - 31.2)^2}{31.2} = \frac{116.64}{31.2} = 3.7385 \\
\frac{(O - E)^2}{E}_{12} &= \frac{(18 - 28.8)^2}{28.8} = \frac{116.64}{28.8} = 4.05 \\
\frac{(O - E)^2}{E}_{21} &= \frac{(10 - 20.8)^2}{20.8} = \frac{116.64}{20.8} = 5.6077 \\
\frac{(O - E)^2}{E}_{22} &= \frac{(30 - 19.2)^2}{19.2} = \frac{116.64}{19.2} = 6.075 \\
\chi^2 &= 3.7385 + 4.05 + 5.6077 + 6.075 \\
 &= 19.471 \\
\mathrm{df} &= (2 - 1)(2 - 1) = 1 \\
p &= P(\chi^2_1 > 19.471) = 1.02\times 10^{-5} \\
V &= \sqrt{\frac{\chi^2}{N}} = \sqrt{\frac{19.471}{100}} = 0.441
\end{align*}


\section{Benjamini--Hochberg: many p-values}
\input{generated/nb_by_hand_10}
\subsection*{The working, typeset}
\begin{align*}
m &= 6, \quad q = 0.05 \\
p_{(1)} = 0.001 &= \frac{1}{6} \times 0.05 = 0.0083 \qquad \text{\small (pass)} \\
p_{(2)} = 0.008 &= \frac{2}{6} \times 0.05 = 0.0167 \qquad \text{\small (pass)} \\
p_{(3)} = 0.039 &= \frac{3}{6} \times 0.05 = 0.025 \qquad \text{\small (fail)} \\
p_{(4)} = 0.041 &= \frac{4}{6} \times 0.05 = 0.0333 \qquad \text{\small (fail)} \\
p_{(5)} = 0.042 &= \frac{5}{6} \times 0.05 = 0.0417 \qquad \text{\small (fail)} \\
p_{(6)} = 0.6 &= \frac{6}{6} \times 0.05 = 0.05 \qquad \text{\small (fail)} \\
k_{\max} &= 2 \qquad \text{\small (keep the 2 smallest p-values)}
\end{align*}

